\documentclass[12pt, a4paper]{article}

% ─── Packages ───────────────────────────────────────────────────────────────
\usepackage[utf8]{inputenc}
\usepackage[T1]{fontenc}
\usepackage{lmodern}
\usepackage[a4paper, margin=2.5cm]{geometry}
\usepackage{amsmath, amssymb, amsthm}
\usepackage{graphicx}
\usepackage{float}
\usepackage{booktabs}
\usepackage{array}
\usepackage{hyperref}
\usepackage{cite}
\usepackage{xcolor}
\usepackage{algorithm}
\usepackage{algpseudocode}
\usepackage{caption}
\usepackage{subcaption}
\usepackage{setspace}
\usepackage{titlesec}
\usepackage{fancyhdr}
\usepackage{enumitem}
\usepackage{tikz}

% ─── Page Style ─────────────────────────────────────────────────────────────
\pagestyle{fancy}
\fancyhf{}
\rhead{\small IRS Beamforming Optimization}
\lhead{\small [Your Name] -- [Student ID]}
\cfoot{\thepage}

\onehalfspacing

% ─── Theorem Environments ───────────────────────────────────────────────────
\newtheorem{proposition}{Proposition}
\newtheorem{remark}{Remark}

% ─── Custom Commands ────────────────────────────────────────────────────────
\newcommand{\vect}[1]{\mathbf{#1}}          % bold vectors
\newcommand{\mat}[1]{\mathbf{#1}}           % bold matrices
\newcommand{\norm}[1]{\left\|#1\right\|}    % norm
\newcommand{\abs}[1]{\left|#1\right|}       % absolute value
\newcommand{\TODO}[1]{\textcolor{red}{\textbf{[TODO: #1]}}}  % placeholder marker

% ════════════════════════════════════════════════════════════════════════════
\begin{document}

% ─── Title Page ─────────────────────────────────────────────────────────────
\begin{titlepage}
    \centering
    \vspace*{1cm}

    {\Large \textbf{[Your University Name]}}\\[0.3cm]
    {\large Department of Electrical Engineering / Electronics \& Telecommunications}\\[2cm]

    \rule{\linewidth}{0.5mm}\\[0.4cm]
    {\LARGE \textbf{Applications of Electronic Circuits\\[0.3cm]
    in Communication Engineering}}\\[0.3cm]
    \rule{\linewidth}{0.5mm}\\[1.5cm]

    {\large \textbf{Project Title:}}\\[0.3cm]
    {\Large Beamforming Optimization for IRS-Aided Wireless Systems\\
    Using Particle Swarm Optimization}\\[2cm]   % <-- change algorithm name if needed

    \begin{tabular}{rl}
        \textbf{Student Name:}  & \TODO{Your Full Name} \\[0.2cm]
        \textbf{Student ID:}    & \TODO{Your Student ID} \\[0.2cm]
        \textbf{Class:}         & \TODO{Your Class Code} \\[0.2cm]
        \textbf{Instructor:}    & \TODO{Instructor Name} \\[0.2cm]
        \textbf{Date:}          & \today \\
    \end{tabular}

    \vfill
    {\small Submitted in partial fulfillment of the requirements for\\
    the course on Applications of Electronic Circuits in Communication Engineering}
\end{titlepage}

% ─── Abstract ───────────────────────────────────────────────────────────────
\begin{abstract}
This report presents the study and implementation of beamforming optimization
for an Intelligent Reflecting Surface (IRS) aided wireless communication system.
We first review the practical phase shift model proposed by Abeywickrama \textit{et al.}
\cite{abeywickrama2020}, which captures the amplitude-phase coupling in IRS reflecting
elements -- a critical hardware imperfection ignored by most prior work.
We then apply \textbf{Particle Swarm Optimization (PSO)}
% ^^^ CHANGE THIS to your chosen algorithm
to jointly optimize the transmit beamforming at the access point and the
IRS phase shifts, with the goal of maximizing the achievable spectral efficiency.
Simulation results are compared against the Alternating Optimization (AO)
baseline reported in the original paper.
Our results show that \TODO{summarize your key finding here after running simulations}.
\end{abstract}

\tableofcontents
\newpage

% ════════════════════════════════════════════════════════════════════════════
\section{Introduction}
\label{sec:intro}

\subsection{Motivation}
Wireless communication systems constantly demand higher data rates, better coverage,
and improved energy efficiency. Traditional solutions -- such as deploying more base
stations or using relay nodes -- are costly and power-hungry. \textit{Intelligent
Reflecting Surfaces} (IRS) have recently emerged as a low-cost, passive alternative
that can reshape the wireless propagation environment without requiring active
radio-frequency components.

An IRS consists of a large number of passive reflecting elements, each capable of
independently adjusting the phase (and to some extent the amplitude) of the reflected
signal. By coordinating these elements, an IRS can steer the reflected signals toward
a desired receiver, effectively creating a constructive interference channel even in
heavily obstructed environments.

\subsection{Problem Statement}
Most existing IRS research assumes an \textit{ideal phase shift model}, where each
element reflects with unity amplitude regardless of its phase shift. However, as
demonstrated in \cite{abeywickrama2020}, practical IRS elements exhibit a strong
coupling between amplitude and phase shift due to inherent hardware losses. Designing
beamforming strategies under this ideal assumption leads to measurable performance
degradation in real systems.

This project addresses the joint optimization of:
\begin{enumerate}
    \item \textbf{Transmit beamforming} at the access point (AP), and
    \item \textbf{IRS phase shifts}, under the practical phase shift model,
\end{enumerate}
to maximize the achievable rate at a single-antenna user.

\subsection{Chosen Algorithm}
To solve this non-convex optimization problem, we employ
\textbf{Particle Swarm Optimization (PSO)}, a population-based metaheuristic
algorithm inspired by the collective behavior of bird flocks.
% ^^^ CHANGE THIS section if you chose a different algorithm

\subsection{Report Organization}
The remainder of this report is organized as follows.
Section~\ref{sec:sysmodel} describes the system model and the practical phase
shift model. Section~\ref{sec:algorithm} presents the chosen optimization algorithm
in detail. Section~\ref{sec:optimization} formulates the optimization problem and
explains how the algorithm is applied. Section~\ref{sec:results} presents and
discusses simulation results. Section~\ref{sec:conclusion} concludes the report.

% ════════════════════════════════════════════════════════════════════════════
\section{System Model}
\label{sec:sysmodel}

\subsection{Network Topology}
We consider a Multiple-Input Single-Output (MISO) downlink system consisting of:
\begin{itemize}
    \item An \textbf{Access Point (AP)} equipped with $M$ antennas,
    \item An \textbf{IRS} composed of $N$ passive reflecting elements,
    \item A \textbf{single-antenna user}.
\end{itemize}

The system geometry is illustrated in Figure~\ref{fig:system_model}. The IRS is
deployed between the AP and user to assist communication, particularly when the
direct AP-user link is weak due to obstacles.

\begin{figure}[H]
    \centering
    % Replace with your own diagram or simulation screenshot
    \TODO{Insert system diagram here -- e.g., exported from Python/draw.io}
    \caption{IRS-aided MISO wireless communication system.}
    \label{fig:system_model}
\end{figure}

\subsection{Channel Model}
Three channels are defined:
\begin{itemize}
    \item $\vect{h}_d \in \mathbb{C}^{M \times 1}$: the direct channel from AP to user,
    \item $\vect{h}_r \in \mathbb{C}^{N \times 1}$: the channel from IRS to user,
    \item $\mat{G} \in \mathbb{C}^{N \times M}$: the channel from AP to IRS.
\end{itemize}
All channels follow a quasi-static flat-fading model and are assumed known at the AP.
Path loss is modeled with a reference attenuation of 40~dB at 1~m, and path loss
exponents of 2.2 (AP-IRS), 2.8 (IRS-user), and 3.8 (AP-user) \cite{abeywickrama2020}.

\subsection{Signal Model}
The AP transmits signal $x = \vect{w}s$, where $\vect{w} \in \mathbb{C}^{M \times 1}$
is the beamforming vector and $s$ is the data symbol with $\mathbb{E}[|s|^2] = 1$.
The received signal at the user is:
\begin{equation}
    y = \left(\vect{v}^H \boldsymbol{\Phi} + \vect{h}_d^H\right)\vect{w}s + z,
    \label{eq:received_signal}
\end{equation}
where $\boldsymbol{\Phi} = \text{diag}(\vect{h}_r^H)\mat{G}$ combines the AP-to-IRS
and IRS-to-user channels, $\vect{v} \in \mathbb{C}^{N \times 1}$ is the IRS
reflection coefficient vector, and $z \sim \mathcal{CN}(0, \sigma^2)$ is
additive white Gaussian noise.

\subsection{Achievable Rate}
The achievable spectral efficiency (bits/s/Hz) is:
\begin{equation}
    R = \log_2\!\left(1 + \frac{\abs{(\vect{v}^H\boldsymbol{\Phi}
        + \vect{h}_d^H)\vect{w}}^2}{\sigma^2}\right).
    \label{eq:rate}
\end{equation}
Maximizing $R$ is equivalent to maximizing the received SNR.

% ─── Practical Phase Shift Model ────────────────────────────────────────────
\subsection{Practical Phase Shift Model}
\label{subsec:psm}

Unlike the ideal model which assumes $|v_n| = 1$ for all phase shifts, the
practical model \cite{abeywickrama2020} captures the hardware reality that
amplitude and phase are coupled. The $n$-th element's reflection coefficient is:
\begin{equation}
    v_n = \beta_n(\theta_n) e^{j\theta_n}, \quad \theta_n \in [-\pi, \pi),
    \label{eq:vn}
\end{equation}
where the amplitude $\beta_n(\theta_n)$ is given by:
\begin{equation}
    \beta_n(\theta_n) = (1 - \beta_{\min})
        \left(\frac{\sin(\theta_n - \phi) + 1}{2}\right)^k + \beta_{\min}.
    \label{eq:beta}
\end{equation}

The parameters $\beta_{\min}$, $\phi$, and $k$ characterize the hardware:
\begin{itemize}
    \item $\beta_{\min} \geq 0$: the minimum reflection amplitude (floor),
    \item $\phi \geq 0$: horizontal offset of the amplitude minimum,
    \item $k \geq 0$: controls the sharpness of the amplitude dip.
        When $k=0$, \eqref{eq:beta} reduces to the ideal model ($\beta_n = 1$).
\end{itemize}

In our simulations, we use $\beta_{\min} = 0.2$, $k = 1.6$, and
$\phi = 0.43\pi$, as reported in \cite{abeywickrama2020}.

\begin{remark}
    The minimum amplitude occurs near $\theta_n = 0$ and the maximum near
    $\theta_n = \pm\pi$. This means simply setting $\theta_n = \angle\varphi_n$
    (as in the ideal design) is suboptimal when $\angle\varphi_n \approx 0$.
\end{remark}

% ════════════════════════════════════════════════════════════════════════════
\section{Optimization Algorithm: Particle Swarm Optimization}
\label{sec:algorithm}
% ^^^ RENAME THIS SECTION to match your chosen algorithm

\subsection{Basic Description}
Particle Swarm Optimization (PSO) \cite{kennedy1995} is a population-based
metaheuristic that simulates the social behavior of a flock of birds searching
for food. Each candidate solution is represented as a \textit{particle} moving
through the $N$-dimensional search space of IRS phase shifts.

Each particle $i$ maintains:
\begin{itemize}
    \item $\vect{x}_i$: its current position (a candidate phase shift vector),
    \item $\vect{p}_i$: its personal best position found so far,
    \item $\vect{g}$: the global best position found by any particle,
    \item $\vect{v}_i^{\text{vel}}$: its velocity (direction of movement).
\end{itemize}

\subsection{Update Rules}
At each iteration $t$, velocities and positions are updated as:
\begin{align}
    \vect{v}_i^{(t+1)} &= w\,\vect{v}_i^{(t)}
        + c_1 r_1 \left(\vect{p}_i - \vect{x}_i^{(t)}\right)
        + c_2 r_2 \left(\vect{g}   - \vect{x}_i^{(t)}\right),
        \label{eq:pso_vel} \\
    \vect{x}_i^{(t+1)} &= \vect{x}_i^{(t)} + \vect{v}_i^{(t+1)},
        \label{eq:pso_pos}
\end{align}
where $w$ is the inertia weight, $c_1$ and $c_2$ are cognitive and social
coefficients, and $r_1, r_2 \sim \mathcal{U}(0,1)$ are random scalars.

\subsection{Pseudocode}
\begin{algorithm}[H]
\caption{PSO for IRS Phase Shift Optimization}
\label{alg:pso}
\begin{algorithmic}[1]
\Require $N$ (IRS elements), $S$ (swarm size), $T$ (max iterations),
         channels $\vect{h}_d, \vect{h}_r, \mat{G}$,
         model params $\beta_{\min}, \phi, k$
\Ensure  Optimal phase shifts $\vect{\theta}^*$, maximum rate $R^*$

\State \textbf{Initialize:} for each particle $i = 1,\ldots,S$:
    $\vect{x}_i \leftarrow \text{Uniform}(-\pi, \pi)^N$,
    $\vect{v}_i^{\text{vel}} \leftarrow \vect{0}$,
    $\vect{p}_i \leftarrow \vect{x}_i$
\State $\vect{g} \leftarrow \arg\max_i f(\vect{x}_i)$

\For{$t = 1$ \textbf{to} $T$}
    \For{$i = 1$ \textbf{to} $S$}
        \State Compute $\beta_n$ from \eqref{eq:beta} for each element $n$
        \State Build reflection vector: $v_n = \beta_n e^{j x_{i,n}}$
        \State Compute combined channel:
               $\vect{c} = \vect{v}^H\boldsymbol{\Phi} + \vect{h}_d^H$
        \State Compute MRT beamforming:
               $\vect{w}^* = \sqrt{P_T}\,\vect{c}^H / \norm{\vect{c}}$
        \State Evaluate fitness: $f_i = \log_2(1 + |\vect{c}\vect{w}^*|^2/\sigma^2)$
        \If{$f_i > f(\vect{p}_i)$}
            $\vect{p}_i \leftarrow \vect{x}_i$
        \EndIf
        \If{$f_i > f(\vect{g})$}
            $\vect{g} \leftarrow \vect{x}_i$
        \EndIf
        \State Update velocity via \eqref{eq:pso_vel}
        \State Update position via \eqref{eq:pso_pos}
        \State Clamp: $\vect{x}_i \leftarrow \text{clip}(\vect{x}_i, -\pi, \pi)$
    \EndFor
\EndFor
\State \Return $\vect{g}$, $f(\vect{g})$
\end{algorithmic}
\end{algorithm}

\subsection{Computational Complexity}
The per-iteration cost is $\mathcal{O}(S \cdot N \cdot M)$, where $S$ is the
swarm size, $N$ the number of IRS elements, and $M$ the number of AP antennas.
Total complexity over $T$ iterations is $\mathcal{O}(T \cdot S \cdot N \cdot M)$.
For typical values ($T=200$, $S=30$, $N=40$, $M=2$), this is computationally
tractable on a standard laptop.

\subsection{Convergence Behavior}
PSO is not guaranteed to find the global optimum due to the non-convex nature
of the problem. However, empirically it converges within \TODO{X} iterations
for this problem size (see Figure~\ref{fig:convergence}). The inertia weight
$w$ is linearly decayed from 0.9 to 0.4 over iterations to balance exploration
(early) and exploitation (late).

\subsection{Variants and Parameter Settings}
\TODO{Briefly mention 1-2 PSO variants, e.g., LPSO (local topology), APSO
(adaptive PSO), or PSO with constriction factor. Then state the exact
hyperparameters you used.}

The hyperparameters used in this work are summarized in Table~\ref{tab:params}.

\begin{table}[H]
    \centering
    \caption{PSO hyperparameters used in simulation.}
    \label{tab:params}
    \begin{tabular}{lll}
        \toprule
        \textbf{Parameter} & \textbf{Symbol} & \textbf{Value} \\
        \midrule
        Swarm size          & $S$   & 30 \\
        Max iterations      & $T$   & 200 \\
        Inertia weight      & $w$   & 0.9 $\to$ 0.4 (linear decay) \\
        Cognitive coeff.    & $c_1$ & 2.0 \\
        Social coeff.       & $c_2$ & 2.0 \\
        \bottomrule
    \end{tabular}
\end{table}

% ════════════════════════════════════════════════════════════════════════════
\section{Problem Formulation and Algorithm Application}
\label{sec:optimization}

\subsection{Joint Beamforming Optimization Problem}
The joint optimization of AP beamforming $\vect{w}$ and IRS phase shifts
$\{\theta_n\}$ is formulated as:
\begin{equation}
    \max_{\vect{w},\,\vect{v},\,\{\theta_n\}}\;
        \abs{\left(\vect{v}^H\boldsymbol{\Phi} + \vect{h}_d^H\right)\vect{w}}^2
    \quad \text{s.t.} \quad
    \norm{\vect{w}}^2 \leq P_T, \quad
    v_n = \beta_n(\theta_n)e^{j\theta_n},\; \theta_n \in [-\pi,\pi).
    \label{eq:P0}
\end{equation}

\subsection{Decoupling via MRT Beamforming}
For any fixed IRS phase shift vector $\vect{v}$, the optimal transmit
beamforming is Maximum Ratio Transmission (MRT):
\begin{equation}
    \vect{w}^* = \sqrt{P_T}\,
        \frac{\left(\vect{v}^H\boldsymbol{\Phi} + \vect{h}_d^H\right)^H}
             {\norm{\vect{v}^H\boldsymbol{\Phi} + \vect{h}_d^H}}.
    \label{eq:MRT}
\end{equation}
Substituting \eqref{eq:MRT} into \eqref{eq:P0} reduces the problem to
optimizing only the IRS phase shifts $\{\theta_n\}$, which PSO handles directly.

\subsection{PSO Applied to IRS Phase Shift Optimization}
Each PSO particle encodes a candidate phase shift vector
$\vect{x}_i = [\theta_1, \ldots, \theta_N] \in [-\pi, \pi)^N$.
The fitness function evaluated for each particle is the achievable rate
$R$ in \eqref{eq:rate}, computed using the practical model \eqref{eq:beta}
and the MRT beamforming \eqref{eq:MRT}.

PSO searches the $N$-dimensional continuous space for the phase configuration
that maximizes $R$, without requiring gradient information or convexity
assumptions.

% ════════════════════════════════════════════════════════════════════════════
\section{Simulation Results}
\label{sec:results}

\subsection{Simulation Setup}
The simulation parameters follow those of \cite{abeywickrama2020} exactly,
to allow direct comparison:

\begin{table}[H]
    \centering
    \caption{Simulation parameters.}
    \label{tab:sim_params}
    \begin{tabular}{lll}
        \toprule
        \textbf{Parameter} & \textbf{Symbol} & \textbf{Value} \\
        \midrule
        AP antennas              & $M$           & 2 \\
        IRS elements             & $N$           & 40 (Fig.~5), varied (Fig.~6) \\
        AP-IRS distance          &               & 500 m \\
        Reference path loss      &               & 40 dB at 1 m \\
        Path loss exp. (AP-IRS)  & $\alpha_1$    & 2.2 \\
        Path loss exp. (IRS-User)& $\alpha_2$    & 2.8 \\
        Path loss exp. (AP-User) & $\alpha_3$    & 3.8 \\
        Transmit power           & $P_T$         & 36 dBm \\
        Noise variance           & $\sigma^2$    & $-94$ dBm \\
        Phase shift model params & $\beta_{\min}, k, \phi$ & $0.2,\;1.6,\;0.43\pi$ \\
        Channel realizations     &               & 1000 \\
        \bottomrule
    \end{tabular}
\end{table}

\subsection{Achievable Rate vs. AP-User Distance}

\begin{figure}[H]
    \centering
    \TODO{Insert your reproduced Figure 5 here}
    \caption{Achievable rate vs.\ AP-user horizontal distance ($N=40$).
             Comparison of: (1) Ideal IRS upper bound, (2) PSO with practical
             model, (3) AO with practical model (paper baseline), (4) Ideal
             model assumption on practical hardware, (5) No IRS lower bound.}
    \label{fig:rate_vs_distance}
\end{figure}

\TODO{Discuss your Figure~\ref{fig:rate_vs_distance} here.
Compare your PSO curve against the AO curves from the paper.
Comment on: (a) where PSO matches AO, (b) where it diverges and why,
(c) the performance gap between practical and ideal model.}

\subsection{Achievable Rate vs. Number of IRS Elements}

\begin{figure}[H]
    \centering
    \TODO{Insert your reproduced Figure 6 here}
    \caption{Achievable rate vs.\ number of reflecting elements ($d = 498$~m).}
    \label{fig:rate_vs_N}
\end{figure}

\TODO{Discuss Figure~\ref{fig:rate_vs_N}: how does performance scale with $N$?
Does the practical model gap grow with $N$? Why?}

\subsection{Convergence of PSO}

\begin{figure}[H]
    \centering
    \TODO{Insert convergence plot here: best fitness vs. iteration number}
    \caption{PSO convergence behavior for a single channel realization ($N=40$,
             $d=498$~m).}
    \label{fig:convergence}
\end{figure}

\TODO{Comment on the number of iterations needed for convergence and
the effect of swarm size on convergence speed.}

\subsection{Discussion}
\TODO{Write 1-2 paragraphs summarizing:
\begin{itemize}
    \item How well does PSO reproduce the paper's results?
    \item What are PSO's advantages over AO for this problem
          (e.g., no closed-form derivation needed, handles discrete phases)?
    \item What are its disadvantages (e.g., slower convergence, hyperparameter sensitivity)?
\end{itemize}}

% ════════════════════════════════════════════════════════════════════════════
\section{Conclusion}
\label{sec:conclusion}

This report studied the beamforming optimization problem for an IRS-aided MISO
wireless communication system under the practical phase shift model proposed
in \cite{abeywickrama2020}. The key contributions of this work are:
\begin{enumerate}
    \item A summary of the IRS system model and the practical phase shift model,
          highlighting the amplitude-phase coupling that invalidates the ideal assumption.
    \item A detailed description and implementation of Particle Swarm Optimization
          for solving the joint transmit and reflect beamforming problem.
    \item Simulation results that \TODO{describe your finding -- e.g., ``confirm
          that PSO achieves within X\% of the AO baseline while requiring no
          gradient computation''}.
\end{enumerate}

Future work could explore PSO variants with adaptive inertia, hybrid
PSO-AO approaches, or extension to multi-user scenarios.

% ─── References ─────────────────────────────────────────────────────────────
\bibliographystyle{IEEEtran}
\begin{thebibliography}{99}

\bibitem{abeywickrama2020}
S.~Abeywickrama, R.~Zhang, and C.~Yuen,
``Intelligent reflecting surface: Practical phase shift model and beamforming optimization,''
\textit{arXiv:1907.06002v4}, Feb. 2020.

\bibitem{wu2019}
Q.~Wu and R.~Zhang,
``Intelligent reflecting surface enhanced wireless network via joint active and passive beamforming,''
\textit{IEEE Trans. Wireless Commun.}, vol.~18, no.~11, pp.~5394--5409, Nov. 2019.

\bibitem{kennedy1995}
J.~Kennedy and R.~Eberhart,
``Particle swarm optimization,''
in \textit{Proc. IEEE Int. Conf. Neural Networks}, 1995, pp.~1942--1948.

\bibitem{huang2019}
C.~Huang, A.~Zappone, G.~C.~Alexandropoulos, M.~Debbah, and C.~Yuen,
``Reconfigurable intelligent surfaces for energy efficiency in wireless communication,''
\textit{IEEE Trans. Wireless Commun.}, vol.~18, no.~8, pp.~4157--4170, Aug. 2019.

\bibitem{basar2019}
E.~Basar \textit{et al.},
``Wireless communications through reconfigurable intelligent surfaces,''
\textit{IEEE Access}, vol.~7, pp.~116\,753--116\,773, Aug. 2019.

\end{thebibliography}

\end{document}
